\subsection{问题二：全年计划购电策略与紧急购电}

\subsubsection{模型的建立}

问题二去掉日周期约束，把决策期扩展到全年 365 天，并用跨日状态连续性取代它：

\begin{equation}
E_{d,0}=E_{d-1,144},\qquad E_{1/1,0}=6000,\qquad E_{12/31,144}\ \text{自由}
\end{equation}

式 (10) 就是状态转移式在日边界处的一次转移，不需要额外约束。与问题一相比的关键差别在于：储能在不同日期之间可以携带电能，因此“哪天充电、哪天放电”也成为决策的一部分；目标中新增紧急购电项 $\alpha_{em}p_{\tau}q^{em}_{\tau}$，$\alpha_{em}=5$。把跨日时段统一为一个长时域优化问题，是微网与储能调度中处理日间耦合的标准做法\cite{ref-silvente-2015-rolling-horizon,ref-holjevac-2017-receding-horizon,ref-hu-2020-mpc-overview,ref-finnah-2021-da-intraday-adp}。

\subsubsection{紧急购电恒为零的条件与意义}

\begin{quote}
\textbf{定理 2（紧急购电的机制性恒零）} 若 ① 电价 $p_{\tau}>0$；② $\alpha_{em}>1$；③ 计划购电量 $b_{\tau}$ 无上界；④ 计划层对当日负载与光伏为完全信息，则模型存在满足 $q^{em}_{\tau}\equiv0$ 的最优解。

\textbf{证明} 设某时段 $q^{em}_{\tau}=\delta>0$。由于平衡为等式且 $b_{\tau}$ 无上界，把该时段计划购电量改为 $b'_{\tau}=b_{\tau}+\delta$ 仍可行（平衡方程左端不变），而费用变化为 $p_{\tau}\delta-\alpha_{em}p_{\tau}\delta=p_{\tau}\delta(1-\alpha_{em})<0$，即费用严格下降，与最优性矛盾。故存在最优解使 $q^{em}\equiv0$。$\square$
\end{quote}

定理 2 说明：在完全信息且购电无上限的口径下，紧急购电惩罚机制\textbf{不会}被激活，交付表中紧急购电量全零是模型结构的推论，而非数据巧合。这一点必须在论文中解释清楚，否则“表格全零”会被误读为漏填。同时也说明该机制不可删除：一旦引入购电功率上限，紧急购电立即被激活，本文用上限扫描给出激活阈值——限额压到 4218.75 kW 时全年紧急购电量升至 8707.398542 kWh 并涉及 14 天，压到 2000 kW 时升至 7\,845\,369.336456 kWh 并涉及全部 365 天，总费用单调上升至 2.38 倍。因此“无上限”是必须显式声明的前提，而非可以省略的细节。文献亦普遍强调：日前决策与实时结算之间的信息落差与功率约束是应急机制被激活的直接原因\cite{ref-chazarra-2017-perfect-information,ref-toubeau-2020-da-real-time,ref-vanderveen-2016-balancing-market-design,ref-herre-2020-imbalance-settlement,ref-joos-2018-integration-costs}。

\subsubsection{主结果}

全年最优购电费为 $13\,758\,182.573724$ 元，其中交付期（2025 年 2 月 1 日--12 月 31 日，334 天）为 $12\,233\,050.830708$ 元，1 月预热期为 $1\,525\,131.743016$ 元，三者满足“交付期加预热期等于全期”的加总关系。主要总量指标列于表~\ref{tab:p2main}。

\begin{table}[htbp]
\centering
\small
\caption{问题二的主要结果}
\label{tab:p2main}
\begin{tabular}{lll}
\toprule
指标 & 数值 & 口径与说明 \\
\midrule
总购电费（全期 365 天） & 13\,758\,182.573724 元 & 交付期 334 天为 12\,233\,050.830708 元，1 月预热期 1\,525\,131.743016 元 \\
紧急购电费 & 0 元 & 由定理 2 保证 \\
计划购电量 $\sum b$ & 22\,657\,938.157435 kWh & 全期 \\
充电量 $\sum c$ / 放电量 $\sum q^{dis}$ & 7\,364\,620.676347 / 5\,969\,662.747841 kWh & 往返损耗 19\% \\
弃光量 $\sum s$ & 990\,168.090312 kWh & 97 天出现弃光，全部落在交付期 \\
净负荷 $N$ & 20\,272\,812.138617（全期）/ 17\,939\,189.494967（交付期）kWh & 与价格无关，用于跨问题校验 \\
储电量端点 & $E_{2/1,0}=7950.000000$，$E_{12/31,144}=1200.000000$ kWh & 跨日滚动、终端自由 \\
\bottomrule
\end{tabular}
\end{table}

量级自检方面，费用落在独立构造的区间 $[7\,525\,691.131,18\,298\,592.366]$ 元内；套利方向正确，充电加权均价 0.520371 元/kWh 明显低于放电加权均价 1.154986 元/kWh。储电量两端（1200 与 10800 kWh）在最优解中均曾活跃，最大购电量出现在功率上限处，峰值购电功率约 10325.33 kW。

\subsubsection{交付结果}

表~\ref{tab:p2a} 给出题面四个指定日期的六个 10 分钟时段购电量、全天购电量与全天购电费；表~\ref{tab:p2b} 给出这四个日期的分块充放电量。四个日期及全年的紧急购电量均为 0，按填报口径在“紧急购电量”工作表中逐日填“—/0”、每日至少一行。需特别提醒：题面表 4 中的 2025 年 3 月 1 日及三个时段是\textbf{填写格式示例，不是本题答案}。

\begin{table}[htbp]
\centering
\small
\caption{问题二：四个指定日期的指定时段购电量与全天合计（单位：kWh；购电费单位：元）}
\label{tab:p2a}
\begin{tabular}{lrrrrrrrr}
\toprule
日期 & 10:00 & 12:00 & 14:00 & 16:00 & 18:00 & 20:00 & 全天购电量 & 全天购电费 \\
\midrule
2025-03-20 & 0.00 & 558.311667 & 0.00 & 507.324167 & 682.353333 & 0.00 & 61\,658.562881 & 37\,975.400853 \\
2025-06-21 & 0.00 & 0.00 & 0.00 & 27.005000 & 385.901667 & 0.00 & 38\,934.072483 & 20\,470.622813 \\
2025-09-23 & 0.00 & 380.021250 & 0.00 & 493.582500 & 752.645000 & 61.911667 & 66\,265.421309 & 41\,351.897551 \\
2025-12-21 & 0.00 & 906.683333 & 0.00 & 795.930417 & 712.633333 & 30.268333 & 93\,770.214940 & 59\,717.364845 \\
\bottomrule
\end{tabular}
\end{table}

\begin{table}[htbp]
\centering
\footnotesize
\caption{问题二：四个指定日期的分块充放电量（单位：kWh；块顺序为 0--4、4--8、8--12、12--16、16--20、20--24 时）}
\label{tab:p2b}
\begin{tabular}{llrrrrrrl}
\toprule
日期 & 量 & 第 1 块 & 第 2 块 & 第 3 块 & 第 4 块 & 第 5 块 & 第 6 块 & 端点储电量 0:00 $\to$ 24:00 \\
\midrule
2025-03-20 & 充电 & 1666.67 & 833.33 & 6291.46 & 4689.68 & 0.00 & 3503.18 & \multirow{2}{*}{8550.00 $\to$ 4352.87} \\
 & 放电 & 0.00 & 5862.62 & 2777.38 & 254.72 & 5728.06 & 2911.94 & \\
2025-06-21 & 充电 & 0.00 & 833.33 & 701.60 & 9965.07 & 0.00 & 8166.67 & \multirow{2}{*}{2361.06 $\to$ 8550.00} \\
 & 放电 & 0.00 & 1719.95 & 0.00 & 0.00 & 6154.10 & 2485.90 & \\
2025-09-23 & 充电 & 1666.67 & 833.33 & 5245.38 & 5923.78 & 0.00 & 8166.67 & \multirow{2}{*}{8550.00 $\to$ 8550.00} \\
 & 放电 & 0.00 & 6746.77 & 1896.62 & 403.62 & 5640.00 & 3000.00 & \\
2025-12-21 & 充电 & 1666.67 & 1666.67 & 5833.33 & 7574.21 & 0.00 & 8166.67 & \multirow{2}{*}{8550.00 $\to$ 8550.00} \\
 & 放电 & 0.00 & 1425.00 & 7890.00 & 2220.11 & 5797.47 & 2842.53 & \\
\bottomrule
\end{tabular}
\end{table}

这里必须区分两种端点语义：问题一的 0:00 与 24:00 是\textbf{日周期约束}的端点，恒为 6000 kWh；问题二的端点是全年滚动轨迹在某一日的切片，随日期变化，两者不可直接比较。

\subsubsection{策略形态与机制解释}

问题二的最优策略在更长时间尺度上重复了套利逻辑，但也出现两个新现象（图~\ref{fig:p2a}）。

第一，\textbf{储能存在明显的“回归水平”}：334 个交付日中，日初或日末储电量等于 8550.00 kWh 的各有 219 天，全年日初储电量只取 96 个不同水平。这说明最优策略倾向于把储能维持在高位以待高价时段，而不是每天放空再充满；只有 12 月 31 日（决策期最后一天）日末落到下限 1200 kWh，即期末放空，对应式 (8) 中 $-4320$ kWh 的库存变现项。

第二，\textbf{弃光只发生在少数日子}：全期仅 97 天出现弃光，合计 990168.090312 kWh，占光伏总发电量的比例很小。这说明在无售电收益的设定下储能容量基本足以吸收盈余；但弃光变量仍必须保留，否则那些盈余超过储能吸收能力的时段会使模型不可行。

策略形态的统计量为：电价与充电量的相关系数 $-0.3835$，与放电量的相关系数 $+0.6132$；最便宜的 20\% 电价时段承担 48.12\% 的充电量，最贵的 20\% 时段承担 68.42\% 的放电量；交付期单日费用最小 11\,900.552329 元、最大 60\,014.368886 元、均值 36\,625.900691 元（合计与交付期费用逐位一致）。

\begin{figure}[htbp]
\centering
\includegraphics[width=0.80\textwidth]{fig06_p2_energy_flow.png}
\caption{全年逐日供给构成、负载与逐日购电费}
\label{fig:p2a}
\end{figure}

\begin{figure}[htbp]
\centering
\includegraphics[width=0.80\textwidth]{fig08_p2_spec_dates.png}
\caption{题面四个指定日期的日内功率构成}
\label{fig:p2c}
\end{figure}

\subsubsection{模型族与结构对照}

本文设置了 17 组对照（表~\ref{tab:p2ab}），结论有三点。

其一，\textbf{核心结论不依赖模型族与实现路线}。用不同变量排序、消去弃光变量、改用对偶单纯形法的独立第二实现给出的费用与主模型逐位相同；把全年问题精确分解为 365 个单日问题（耦合量取全年最优值）也给出相同结果。差异为零说明主结果不是实现假象。

其二，\textbf{跨日与终端口径有可量化的经济价值}。逐日解耦并允许日终自由递推会使费用上升 26\,894.74 元，这是“跨日携带电能”的价值；逐日施加日周期约束会使费用上升 17\,135.02 元，这是跨日衔接与逐日终端两项的\emph{合并}贡献；把全年终端固定为 6000 kWh 会使费用上升 2217.08 元，对应期末放空的 4320 kWh 库存。三者对应不同的约束改动，不可互相替代或相加。

其三，\textbf{购电上限是唯一能激活紧急购电的机制}（图~\ref{fig:p2d}）。上限在 10326 kW 及以上时不产生任何影响（该值只是“非绑定阈值”），而在 4218.75 kW 及以下时紧急购电被激活且费用随上限下降单调上升，代价从 $+694\,101$ 元迅速放大到 $+19\,011\,429$ 元（上限 2000 kW）。这说明“紧急购电为零”的结论对上限极度敏感，必须与前提条件一并陈述。

\begin{table}[htbp]
\centering
\small
\caption{问题二的时域组织、终端条件与购电上限对照（费用为全期，单位：元）}
\label{tab:p2ab}
\begin{tabular}{p{7.0cm}rr}
\toprule
对照项 & 费用 & 相对基准 \\
\midrule
基准（全年单一长时域 LP） & 13\,758\,182.573724 & --- \\
独立第二实现 / 精确逐日分解 & 13\,758\,182.573724 & $+0.000000\%$ \\
逐日解耦、日终自由递推 & 13\,785\,077.313742 & $+0.195482\%$ \\
逐日施加日周期约束 / 终端固定为 6000 kWh & 13\,775\,317.596266 / 13\,760\,399.657058 & $+0.124544\%$ / $+0.016115\%$ \\
购电上限 4375 / 4218.75 / 2000 kW & 14\,452\,283.680247 / 14\,647\,566.902931 / 32\,769\,611.572155 & $+5.045\%$ / $+6.464\%$ / $+138.183\%$ \\
互补约束整数模型（52\,560 个二元变量） & 13\,758\,186.928903 & $+0.000032\%$ \\
\bottomrule
\end{tabular}
\end{table}

\begin{figure}[htbp]
\centering
\includegraphics[width=0.80\textwidth]{fig09_p2_bcap.png}
\caption{外网购电上限对紧急购电量与总费用的影响：上限降至 4218.75 kW 及以下时紧急购电被激活}
\label{fig:p2d}
\end{figure}

\subsubsection{稳健性与适用边界}

149 次全年线性规划的扰动实验得到两点结论（判据与阈值在跑数前冻结）。第一，\textbf{年度与交付期总量稳健}：参数扰动下费用最大偏移 14.078\%（对应 $\eta=0.72$ 的极端档），噪声情景下费用区间半宽不超过 0.96\%，全部情景可行、无 NaN 或无穷值。第二，\textbf{日尺度交付量对效率口径敏感}：$\eta$ 变化 $\pm10\%$ 可使单日购电量与购电费移动约 $\pm20\%$（实测最大 19.62\%），而年度总量仅移动约 $\pm7\%$。因此表~\ref{tab:p2a}、表~\ref{tab:p2b} 中的单日数值须与该敏感性一并引用。

本文如实登记两条未通过的判据。其一，可行率判据的唯一失败样本来自求解器设置对照，其目标值与基线\textbf{逐位相同}、约束残差为 $2\times10^{-12}$、解可行，唯一未通过的检查是“返回解是否满足充放电互补性”——由于引理 1 只断言\emph{存在}互补最优解而非任意返回解都互补，这属于退化最优面上的返回解选择问题，不构成模型不可行。其二，交付表稳定判据的失败样本来自 $\eta=0.99$ 的情景，该情景使 2025 年 6 月 21 日全天购电量下降 19.10\%，属\emph{有利方向}的变化，但因判据取绝对值口径而被判失败；其反向（$\eta=0.81$）最多使单日购电费上升 12.97\%，有 21 天升幅超过 10\%。据此本问的稳定性分级为“脆弱（需缩小适用范围）”，判据未做事后修改。

此外登记两项属于结构性口径的发现：一是当允许以 0.50 元/kWh 售电时，松散的弃光上界会出现“售出电量超过物理盈余”的非物理套利（售出 5\,845\,778.709901 kWh 远超物理盈余 3\,120\,221.003233 kWh），须先收紧上界，该情景收益不得当作真实售电收入；二是当把紧急购电倍数改为 0.5 时全年紧急购电量恰等于基准的全年计划购电量，这是同一物理计划被整体改记为紧急购电的“重标定效应”，不得把它解读为净负荷。
